\section{Extensión de la gramática}\label{sec:extension}

\subsection{Tipos de expresión}
La versión añade constructores conservando los usos anteriores. Esta clasificación
es normativa para escribir; LaTeX dibuja expresiones y no verifica sus tipos.
\begin{longtable}{@{}p{29mm}>{\raggedright\arraybackslash}p{\dimexpr\linewidth-29mm-2\tabcolsep\relax}@{}}
\toprule Tipo & Construcciones y usos \\ \midrule\endhead
Término & Referente, variable, dato convencional, concepto, propiedad o colección.
Los subtipos importan: una propiedad no es una proposición y una lista no es un conjunto. \\[4pt]
Proposición & Relación aplicada a términos; fórmula convencional; negación,
conjunción, condicional, alternativa, cuantificación y posibilidad bien formada. \\[4pt]
Programa & Asignación, secuencia y mientras. Solo tienen la interpretación de
ejecución descrita en este capítulo. \\[4pt]
Unidad superior & Proposición afirmada, término mostrado, programa escrito,
definición, atribución o pregunta. Se separan por párrafos explícitos. \\[4pt]
\bottomrule
\end{longtable}
Las referencias deben existir en la leyenda. Los argumentos deben respetar la
cantidad, el orden y el tipo del catálogo. Un cuantificador recorre un conjunto;
su cuerpo es una proposición. No se puede pasar un programa a K11 ni una pregunta
a K12. Las atribuciones K16, preguntas K19 y definiciones K09 permanecen en el
nivel superior. Una fuente identificada puede atribuir una proposición modal;
el modelo correspondiente también debe estar declarado.

\subsection{Gramática abstracta de los añadidos}
En este esquema los nombres son metanotación. \(t\) es término, \(D\) un
conjunto, \(\varphi\) una proposición, \(\alpha\) un programa y \(x\) variable.
La agrupación escrita es siempre explícita; no hay precedencia implícita entre
operadores prefijos.
\begin{align*}
D &::= C_i\mid\operatorname{conjunto}(t_1,\ldots,t_n)
       \mid\operatorname{seleccion}(x,D,\varphi)\mid D_{\rm conv}\\
\varphi &::= R_i(t_1,\ldots,t_n)\mid\varphi_{\rm conv}
       \mid\neg\varphi\mid(\varphi\land\varphi)\mid(\varphi\to\varphi)\\
 &\phantom{::=}\mid(\varphi\lor\varphi)\mid\forall x\in D:\varphi
       \mid\exists x\in D:\varphi\mid\Diamond_m\varphi\\
\alpha &::= x:=e\mid(\alpha;\alpha)\mid\operatorname{mientras}(b,\alpha).
\end{align*}
El caso modal tiene las restricciones que siguen. El booleano \(b\) de un programa
es una combinación finita de comparaciones enteras con negación, conjunción o
alternativa; no incluye preguntas, fuentes, cuantificadores ni modalidades.
Los términos enteros \(e\) usan constantes, variables inicializadas, suma, resta
y multiplicación. Se excluyen desbordamiento, división y efectos externos.

\subsection{Alternativas y preguntas}
K17 acepta dos proposiciones y expresa disyunción inclusiva. Es falsa únicamente
cuando ambas son falsas. Su expansión clásica es
\(P\lor Q\equiv\neg(\neg P\land\neg Q)\).
La alternativa exclusiva se compone como
\((P\lor Q)\land\neg(P\land Q)\); no requiere otro signo.

K19 toma una proposición y solicita decidirla. Una respuesta afirmativa confirma
el contenido y una negativa confirma su negación. La ausencia de respuesta o
información no convierte la proposición en falsa. Esta versión no define todavía
preguntas de «quién», «cuánto» o respuestas parciales con un signo propio.

\subsection{Posibilidad con un modelo explícito}
K18 lleva un índice \(m\) de modelo modal. El contexto especifica estados \(W\),
estado actual \(w\), relación de accesibilidad \(A_m\) y valoración \(V\) de los
átomos en cada estado. Su interpretación es:
\[
 w\models\Diamond_m P
 \quad\Longleftrightarrow\quad
 \exists v\in W:\ (w,v)\in A_m\ \text{y}\ v\models P.
\]
Solo se admite el fragmento proposicional: los átomos carecen de variables libres,
y no se combinan cuantificadores con modalidades en una misma fórmula de 0.2.
Todas las modalidades de una fórmula usan el mismo modelo; pueden anidarse.
Si no hay estados accesibles, la posibilidad es falsa. La necesidad puede
abreviarse con \(\neg\Diamond_m\neg P\), sin signo nuevo.

El modelo determina qué significa «posible» en ese documento. No equivale a
probable ni a demostrado. T13 da un modelo finito completamente especificado.
Cambiar el modelo puede cambiar el resultado.

\subsection{Ejecución y cambios de estado}
Un estado asigna valores enteros a las variables del programa. La asignación
actualiza una variable después de evaluar su expresión en el estado anterior.
K21 ejecuta dos programas de izquierda a derecha. K22 comprueba el booleano
antes del cuerpo, ejecuta el cuerpo si es verdadero y vuelve a comprobarlo.
Si es falso al principio, el estado no cambia. Un programa puede no terminar;
el segundo miembro de una secuencia no se ejecuta si el primero no termina.

El condicional K13 conserva la implicación material entre proposiciones. No se
reinterpreta como una instrucción de bifurcación. La repetición con comprobación
final se deriva como \(\alpha;\operatorname{mientras}(\neg b,\alpha)\): ejecuta
el cuerpo al menos una vez. Para ejecutarla a mano, deben conocerse los valores
iniciales de toda variable leída.

T15 termina en \(x=3\), tras tres incrementos. Una medida de progreso es
\(3-x\) mientras \(x<3\); en esa traza es un entero positivo que disminuye.
Para un límite variable \(n\geq0\), inicializar \(x=0\) y aumentar hasta \(n\)
realiza \(n\) incrementos y \(n+1\) comprobaciones. Bajo coste unitario de estas
operaciones, el tiempo es \(\Theta(n+1)\); con enteros arbitrariamente grandes
debe declararse un modelo de coste de bits diferente.

\subsection{Criterios de aceptación y casos inválidos}
\begin{itemize}
\item R2 necesita tres argumentos, con papeles agente, dirección y dato.
\item R5 acepta un conjunto en su segundo lugar; una lista necesita otra operación.
\item Una selección necesita un dominio de conjunto y un cuerpo proposicional.
\item Una variable libre no declarada deja la expresión incompleta.
\item Una afirmación causal R6 necesita el modelo; R1 no permite deducir R6.
\item Una modalidad sin modelo o dentro de cuantificación queda fuera del fragmento.
\item Una asignación no declara una verdad matemática; modifica un estado.
\end{itemize}
